% explainer-source-hash: bce9ca386d29403f9f7ef5635ce05a37dbb97c767bc5d139b36782533ef55271
% explainer-source-commit: ec75488bce3db41a4f254c8c70bbc4f8f1693cf8
% explainer-generated: 2026-07-16
% Regenerate with /explain-all (see WIRING.md). Do not edit this stamp by hand:
% it is how the toolchain knows whether this document still matches the code.
\documentclass[11pt,a4paper]{article}
\input{preamble}

\title{\vspace{-1.5cm}\textbf{University Enrollment System}\\[0.3em]
\large A Deep Dive into a PyQt6 + SQLAlchemy Course Registration Desktop App}
\author{Explainer document for Salman Adnan}
\date{}

\begin{document}
\maketitle
\thispagestyle{empty}

\begin{keyidea}{Read this first: what this document is}
This is an exhaustive, beginner-safe walkthrough of the \file{university-enrollment-system}
project. It assumes no prior knowledge: every technical term is defined at first use.

It is a working document for your own understanding and interview defence, not marketing
copy. Where the code is weak, wrong, or disagrees with the README, this document says so
plainly, in a red box. A flaw you learn about here is a flaw you can defend in a room. A
flaw an interviewer finds first is not.

Two sections carry the weight of the project and are treated in full depth: the
\textbf{data model} (Section~\ref{sec:datamodel}) and the \textbf{enrollment rules}
(Section~\ref{sec:rules}).
\end{keyidea}

\tableofcontents
\newpage

% =====================================================================
\section{What this project is}

\subsection{The one-sentence version}

A desktop application that lets a university student log in, receive a code by email,
confirm it, and then build a course timetable: browse a catalogue of course sections,
put some in a cart, enroll, drop, swap sections, look at the resulting weekly schedule,
export it to a spreadsheet, mark attendance, and read a log of everything they did.

\subsection{Provenance and scope}

Per the project README, this was the semester project for the Databases course
(CS/CE 355/373) at Habib University, built by a two-person team, Salman Adnan and Shayan
Wasif, in the 4th semester. The README records Salman's areas as the email verification
(login MFA) flow, the attendance module with its email confirmations, the Docker
packaging, and integration plus bug fixing across the screens.

\begin{jargon}{README}
A \file{README} is the plain-text introduction file that sits at the root of a code
repository. By convention it is the first thing a reader opens. It is written by humans
and is not executed, so it can drift out of step with the code. Throughout this document,
where the README and the code disagree, \textbf{the code wins}, and the disagreement is
recorded as a finding.
\end{jargon}

By size this is a mid-weight project: 13 Python modules totalling roughly 2,570 lines,
11 interface-definition files, and one pre-seeded database file. That is small enough to
cover completely, which this document does, and large enough that the structure matters.

\subsection{How the facts in this document were established}

This matters, because it tells you which claims you can defend hard and which you cannot.

\begin{itemize}
  \item \textbf{Read.} Every tracked file was read end to end.
  \item \textbf{Queried.} The shipped \file{student.db} was opened read-only and its
        schema, row counts, and actual data were inspected directly. Every number quoted
        about the database in this document came out of the file, not out of the README.
  \item \textbf{Executed.} The data-access layer (\file{queries.py}) was imported into a
        throwaway virtual environment and run against a \emph{copy} of the database, to
        test the capacity arithmetic and the swap operation. The results are reported in
        Sections~\ref{sec:capacity} and~\ref{sec:swap}. The original database was not
        touched.
  \item \textbf{Executed.} PyQt6 was installed in that same throwaway environment purely
        to settle one question of fact about enum truthiness, in
        Section~\ref{sec:checkstate}.
  \item \textbf{Not executed.} The graphical application itself was never launched: PyQt6
        is not installed on the machine this document was written on, and the app needs a
        desktop. Any claim about on-screen behaviour is therefore a reading of the code,
        and is labelled as such.
  \item \textbf{Not executed.} The Docker image was never built: the Docker daemon was not
        running. Section~\ref{sec:docker} contains one clearly labelled inference as a
        result.
\end{itemize}

% =====================================================================
\section{Orientation: the shape of the whole thing}

Before any detail, here is the entire application in one picture. Every box is a window
the user can be looking at; every arrow is a transition the code actually wires up.

\begin{figure}[H]
\centering
\resizebox{\textwidth}{!}{
\begin{tikzpicture}[node distance=8mm]

  \node[box] (login) {\textbf{login\_page.py}\\\scriptsize LoginForm\\\scriptsize email + password};
  \node[box, right=16mm of login] (verif) {\textbf{verification.py}\\\scriptsize VerificationForm\\\scriptsize 6-digit code};
  \node[box, right=16mm of verif, fill=accent!20] (main) {\textbf{main\_window.py}\\\scriptsize MainStudentWindow\\\scriptsize the hub};

  \node[box, above right=6mm and 22mm of main] (add) {\textbf{add\_courses.py}\\\scriptsize AddCoursesWindow\\\scriptsize the cart};
  \node[box, below=5mm of add] (enroll) {\textbf{enroll.py}\\\scriptsize ViewScheduleWindow\\\scriptsize checkout};
  \node[box, below=5mm of enroll] (drop) {\textbf{drop\_courses.py}\\\scriptsize DropCoursesWindow};
  \node[box, below=5mm of drop] (swap) {\textbf{swap\_courses.py}\\\scriptsize SwapCoursesWindow};
  \node[box, below=5mm of swap] (view) {\textbf{view\_schedule.py}\\\scriptsize ViewScheduleWindow};
  \node[box, below=5mm of view] (att) {\textbf{attendance.py}\\\scriptsize AttendanceForm};
  \node[box, below=5mm of att] (log) {\textbf{activity\_log.py}\\\scriptsize ActivityLogWindow};

  \node[store, below=20mm of verif] (q) {\textbf{queries.py}\\\scriptsize AllQueries\\\scriptsize the only way to the data};
  \node[store, below=8mm of q] (db) {\textbf{student.db}\\\scriptsize SQLite};
  \node[extbox, below=20mm of login] (smtp) {\scriptsize Gmail SMTP\\\scriptsize (optional)};

  \draw[flow] (login) -- node[lbl,above]{password\\matches} (verif);
  \draw[flow] (verif) -- node[lbl,above]{code\\matches} (main);
  \draw[dflow] (verif) -- node[lbl,left,pos=0.4]{send code} (smtp);
  \draw[dflow] (att.east) to[out=0,in=0,looseness=1.4] node[lbl,pos=0.5]{confirmations} ($(smtp.east)+(0,0)$);

  \foreach \t in {add,enroll,drop,swap,view,att,log}
    \draw[flow] (main.east) -- (\t.west);

  \draw[flow] (add.south east) to[out=-30,in=30,looseness=1.6] node[lbl]{Next} (enroll.north east);

  \foreach \t in {add,enroll,drop,swap,view,att,log,main,login,verif}
    \draw[dflow, opacity=0.35] (\t) -- (q);

  \draw[flow] (q) -- node[lbl,right]{SQLAlchemy ORM} (db);

\end{tikzpicture}}
\caption{The whole application. Solid arrows are window transitions; faint dashed arrows
are data-access calls. Every screen without exception reaches the database through
\code{AllQueries}. That single choke point is the best structural decision in the project.}
\label{fig:arch}
\end{figure}

\begin{jargon}{Desktop app, GUI, and PyQt6}
A \textbf{GUI} (Graphical User Interface) is software you click on, as opposed to one you
type commands at. \textbf{Qt} is a long-established C++ toolkit for building GUIs: it
supplies the buttons, tables, and windows. \textbf{PyQt6} is a Python binding for Qt
version 6, meaning a translation layer that lets Python code drive Qt's C++ machinery.
This project is a \textbf{desktop app}: it runs as a program on your own computer and
draws its own windows, unlike a web app, which runs in a browser.
\end{jargon}

\begin{jargon}{SQLite}
\textbf{SQLite} is a database that lives in a single ordinary file on disk, with no server
process to install or start. Here that file is \file{student.db}. This makes it ideal for
coursework and for shipping a demo, since a clone of the repository comes with its data
already in it, and less appropriate for anything with many simultaneous users.
\end{jargon}

\begin{jargon}{ORM and SQLAlchemy}
A database stores rows in tables. Python works in objects. An \textbf{ORM}
(Object-Relational Mapper) is the layer that translates between the two, so that a row in
the \code{student} table appears in Python as a \code{Student} object with a \code{.name}
attribute. \textbf{SQLAlchemy} is the standard Python ORM. It means you rarely write SQL
by hand; you write \code{session.query(Student).filter\_by(std\_id=123).first()} and
SQLAlchemy emits the \code{SELECT} for you.
\end{jargon}

\subsection{The file layout}

\begin{figure}[H]
\centering
\begin{forest} dirtree
[university-enrollment-system/
  [login\_page.py \quad{\rmfamily\itshape entry point}]
  [verification.py \quad{\rmfamily\itshape MFA code screen}]
  [main\_window.py \quad{\rmfamily\itshape hub}]
  [add\_courses.py \quad{\rmfamily\itshape the cart}]
  [enroll.py \quad{\rmfamily\itshape checkout}]
  [drop\_courses.py]
  [swap\_courses.py]
  [view\_schedule.py \quad{\rmfamily\itshape + Excel export}]
  [attendance.py]
  [activity\_log.py]
  [queries.py \quad{\rmfamily\itshape data-access layer}]
  [student\_database.py \quad{\rmfamily\itshape ORM models + engine}]
  [courses\_database.py \quad{\rmfamily\itshape DEAD, never imported}]
  [*.ui \quad{\rmfamily\itshape 11 Qt Designer layouts}]
  [student.db \quad{\rmfamily\itshape pre-seeded SQLite file}]
  [requirements.txt]
  [Dockerfile]
  [{run-linux.sh, run-macos.sh, run-windows.ps1}]
  [.env.example]
  [Proposal.pdf]
  [docs/
    [screenshots/app.png]
    [explainers/ \quad{\rmfamily\itshape this document}]
  ]
]
\end{forest}
\caption{Every tracked file. The project is flat: no packages, no \file{src/} directory,
one module per screen.}
\end{figure}

\begin{honest}{The flat layout is a real constraint, not just a style choice}
Every screen calls \code{uic.loadUi('name.ui', self)} with a \emph{relative} path. A
relative path is resolved against the process's current working directory, not against the
location of the Python file. The practical consequence: the application only runs if you
launch it from inside the project directory. Run \code{python /some/where/login\_page.py}
from your home directory and it crashes on the first \code{loadUi}. The README states this
plainly and it is why the Dockerfile sets \code{WORKDIR /app}. It is worth being able to
say out loud: this is a bug that packaging has to work around rather than a decision.
\end{honest}

\begin{jargon}{Qt Designer and \file{.ui} files}
Rather than writing code to place every button, you can draw the window in a visual editor
called \textbf{Qt Designer}, which saves the layout as a \file{.ui} file (an XML document
listing widgets, their names, and their positions). \code{uic.loadUi} reads that file
\emph{at runtime} and builds the widgets, attaching each one to the window object under
its Designer name. That is why \code{self.login\_button} exists in
\file{login\_page.py} even though no line of Python ever created it.
\end{jargon}

% =====================================================================
\section{The data model}
\label{sec:datamodel}

This is a Databases course project, so the schema is the centre of gravity. Everything in
this section was read out of \file{student\_database.py} and then verified against the
actual \file{student.db} file.

\begin{jargon}{Schema, table, row, column}
A \textbf{table} is a grid of data, like a spreadsheet tab. A \textbf{column} is one named
field with a type (a whole number, a string of text, a time). A \textbf{row} is one record.
The \textbf{schema} is the full description of which tables exist and what columns they
have. In SQLAlchemy the schema is written as Python classes, and
\code{Base.metadata.create\_all(engine)} turns those classes into real tables.
\end{jargon}

\begin{jargon}{Primary key and foreign key}
A \textbf{primary key} (PK) is the column, or the combination of columns, that uniquely
identifies a row. No two rows may share it. When a PK is made of more than one column it
is called a \textbf{composite} primary key. A \textbf{foreign key} (FK) is a column that
points at another table's primary key, expressing a relationship: ``this enrollment row
belongs to that student''. A database can \emph{enforce} foreign keys, refusing to store a
row that points at a student who does not exist. Whether this one does is the subject of
Section~\ref{sec:fkbroken}, and the answer will surprise you.
\end{jargon}

\subsection{The five tables, drawn}

\begin{figure}[H]
\centering
\resizebox{\textwidth}{!}{
\begin{tikzpicture}[
  tbl/.style={rectangle, draw=inkblue, fill=white, rounded corners=2pt,
              align=left, inner sep=0pt, font=\scriptsize\ttfamily},
  hdr/.style={font=\scriptsize\bfseries\rmfamily, text=white, fill=inkblue,
              align=center, inner sep=3pt, minimum width=42mm},
  dead/.style={hdr, fill=warnred},
  body/.style={align=left, inner sep=4pt, font=\scriptsize\ttfamily, text width=42mm}
]

% ---- STUDENT
\node[hdr] (sh) {student};
\node[tbl, below=0pt of sh, anchor=north] (sb)
 {\begin{tabular}{@{}p{40mm}@{}}
   \textbf{std\_id} \hfill INTEGER \textit{PK}\\
   name \hfill VARCHAR\\
   email \hfill VARCHAR\\
   password \hfill VARCHAR \textcolor{warnred}{plaintext}\\
   dob \hfill DATE\\
   major \hfill VARCHAR\\[2pt]
   \rmfamily\itshape\scriptsize 6 rows in student.db
 \end{tabular}};

% ---- COURSE_LIST
\node[hdr, right=52mm of sh] (ch) {course\_list \rmfamily\normalfont{\itshape(a section)}};
\node[tbl, below=0pt of ch, anchor=north] (cb)
 {\begin{tabular}{@{}p{40mm}@{}}
   \textbf{course\_list\_id} \hfill INTEGER \textit{PK}\\
   \textbf{section} \hfill VARCHAR \textit{PK}\\
   course\_name \hfill VARCHAR\\
   start\_time \hfill TIME\\
   end\_time \hfill TIME\\
   days \hfill VARCHAR \textcolor{accent}{"MoWe"}\\
   room \hfill VARCHAR\\
   instructor \hfill VARCHAR\\
   status \hfill BOOLEAN\\
   max\_capacity \hfill \textcolor{warnred}{never read}\\
   current\_capacity \hfill INTEGER\\[2pt]
   \rmfamily\itshape\scriptsize 17 rows in student.db
 \end{tabular}};

% ---- COURSE (junction)
\node[hdr, below=26mm of sb.south, anchor=north, fill=accent] (jh) {course \rmfamily\normalfont{\itshape(an enrollment)}};
\node[tbl, below=0pt of jh, anchor=north, draw=accent] (jb)
 {\begin{tabular}{@{}p{40mm}@{}}
   \textbf{course\_id} \hfill INTEGER \textit{PK,FK}\\
   \textbf{std\_id} \hfill INTEGER \textit{PK,FK}\\
   \textbf{section} \hfill VARCHAR \textit{PK}\\
   is\_temp \hfill BOOLEAN \textcolor{accent}{cart flag}\\
   type \hfill VARCHAR \textcolor{accent}{"Add"/NULL}\\
   swap\_id \hfill \textcolor{warnred}{dead, all NULL}\\
   remove\_id \hfill \textcolor{warnred}{dead, all NULL}\\[2pt]
   \rmfamily\itshape\scriptsize 10 rows in student.db
 \end{tabular}};

% ---- ACTIVITY
\node[hdr, right=52mm of jh] (ah) {activity};
\node[tbl, below=0pt of ah, anchor=north] (ab)
 {\begin{tabular}{@{}p{40mm}@{}}
   \textbf{activity\_id} \hfill INTEGER \textit{PK}\\
   std\_id \hfill INTEGER \textit{FK}\\
   course\_id \hfill \textcolor{warnred}{VARCHAR} \textit{FK}\\
   activity\_type \hfill VARCHAR\\
   activity\_date \hfill DATE\\
   activity\_time \hfill TIME\\
   course\_section \hfill VARCHAR\\[2pt]
   \rmfamily\itshape\scriptsize 116 rows in student.db
 \end{tabular}};

% ---- FILTER (dead)
\node[dead, below=14mm of jb.south, anchor=north] (fh) {filter \rmfamily\normalfont{\itshape(DEAD)}};
\node[tbl, below=0pt of fh, anchor=north, draw=warnred] (fb)
 {\begin{tabular}{@{}p{40mm}@{}}
   \textbf{filter\_id} \hfill INTEGER \textit{PK}\\
   std\_id \hfill INTEGER \textit{FK}\\
   by\_day \hfill BOOLEAN\\
   by\_week \hfill BOOLEAN\\[2pt]
   \rmfamily\itshape\scriptsize 0 rows; no code reads
   \rmfamily\itshape\scriptsize or writes this table
 \end{tabular}};

% ---- edges
\draw[flow] (jb.west) to[out=180,in=180,looseness=0.9]
  node[lbl,pos=0.5]{many\\enrollments\\per student} (sb.west);
\draw[flow] (jb.east) -- node[lbl,pos=0.5,above]{many enrollments\\per section} (cb.south west);
\draw[flow] (ab.north) -- node[lbl,pos=0.5,right]{logged for} (cb.south);
\draw[flow] (ab.west) -- node[lbl,pos=0.5,above]{done by} (jb.east);
\draw[dflow, draw=warnred] (fb.west) to[out=180,in=180,looseness=1.7] (sb.west);

\end{tikzpicture}}
\caption{The full schema as it exists in \file{student.db}, with row counts taken from the
shipped file. Blue is the junction table that carries the whole enrollment idea. Red marks
what is declared but dead.}
\label{fig:er}
\end{figure}

\subsection{Table by table}

\subsubsection{\code{student}}

Six columns, one per student, keyed by \code{std\_id}. The shipped database holds 6 rows:
Alice (20180), Bob (123), Steve (153), Babar (173), Umer (257), Ronaldo (273).

\begin{honest}{Passwords are stored in plaintext, and the README says so}
\code{Student.password} is a plain \code{String}. \file{login\_page.py} fetches it with
\code{get\_password\_by\_email(email)} and compares with \code{password == user\_password},
a direct string comparison. There is no hashing, no salt, nothing.

This is genuinely bad, and the right way to talk about it is exactly how the README already
does: it is a coursework limitation, it is called out in the Usage section, and hashing
with bcrypt is listed as a follow-up. Do not defend it as a design choice. Say: ``it is
plaintext, I know why that is wrong, I documented it rather than hiding it, and the fix is
bcrypt at the two call sites.''
\end{honest}

\subsubsection{\code{course\_list}: the catalogue}

This is the table of \emph{sections}, not of courses. The distinction is the single most
important thing to understand about this schema.

\begin{keyidea}{A section, not a course, is the unit of everything}
\code{course\_list} has a \textbf{composite primary key}:
\code{(course\_list\_id, section)}. So ``Data Structure II'' is not one row. It is three:

\begin{center}\ttfamily\small
\begin{tabular}{@{}llllll@{}}
\toprule
id & section & name & days & start & end \\
\midrule
201 & L1 & Data Structure II & MoWe & 11:30 & 12:45 \\
201 & L2 & Data Structure II & WeFr & 14:30 & 15:45 \\
201 & L3 & Data Structure II & MoWe & 16:00 & 17:15 \\
\bottomrule
\end{tabular}
\end{center}

The course \emph{name}, and every timing detail, is repeated on each section row. A student
enrolls in a section. A clash is between sections. A swap, as
Section~\ref{sec:swap} shows, is defined as moving between two sections of the
\emph{same} \code{course\_list\_id}. Once you see that ``a row is a section'', the rest of
the schema reads cleanly.
\end{keyidea}

\begin{honest}{The catalogue is not normalised, and this is a Databases project}
Because the course name lives on every section row, ``Data Structure II'' is stored three
times. Rename the course and you must update three rows or the catalogue disagrees with
itself. The textbook fix is two tables: a \code{course} table
(\code{course\_id}, \code{course\_name}) and a \code{section} table referencing it with
the timing, room, instructor, and capacity.

This is worth rehearsing precisely because of the course it was submitted to. An examiner
looking for normalisation will find this immediately. The honest answer is that the
denormalisation is real, it buys nothing here except simpler queries, and the two-table
split is what you would do now.
\end{honest}

The \code{days} column deserves a note: it is a string of two-character day codes glued
together, so \code{"MoWe"} means Monday and Wednesday. The whole catalogue uses only four
patterns, verified against the file: \code{MoWe}, \code{WeFr}, \code{TuTh}, and one
\code{MoSa}. Section~\ref{sec:clash} shows how the clash checker takes this string apart.

\begin{honest}{Junk in the seeded catalogue}
Row \code{(520, 'S1')} is called ``Test Course'', runs \code{02:30} to \code{02:45} on
\code{MoSa}, and is plainly leftover test data that shipped. It is 15 minutes long, at half
past two in the morning, on Monday and Saturday. If you demo the app live, this row will
appear in the catalogue. Either delete it from the seed or be ready to explain it.
\end{honest}

\subsubsection{\code{course}: the junction table, and the cart}

\begin{jargon}{Junction table (also: join table, associative entity)}
A student takes many sections; a section holds many students. That is a
\textbf{many-to-many} relationship, and a relational database cannot express it with a
single foreign key on either side. The standard solution is a third table whose rows each
pair one student with one section. That is a \textbf{junction table}. Here it is
\code{course}, and its primary key \code{(course\_id, std\_id, section)} is exactly the
statement ``this student holds this section of this course, at most once''.
\end{jargon}

The clever part is that this table does double duty. It holds both what you are
\emph{enrolled in} and what is sitting in your \emph{cart}, distinguished by two columns:

\begin{center}
\begin{tabular}{@{}lll@{}}
\toprule
\textbf{\code{is\_temp}} & \textbf{\code{type}} & \textbf{meaning} \\
\midrule
\code{True} & \code{"Add"} & in the cart, not yet enrolled \\
\code{False} & \code{NULL} & really enrolled \\
\bottomrule
\end{tabular}
\end{center}

Verified against the shipped database: every one of the 10 \code{course} rows has
\code{is\_temp = 0} and \code{type = NULL}. That is, the shipped state has nobody with
anything in their cart, which is what you would expect from a clean save.

\begin{honest}{\code{swap\_id} and \code{remove\_id} are dead columns}
Both are declared on the \code{Course} model, both have foreign keys, and the source even
carries a commented-out earlier version of their definitions. No code anywhere assigns
either one. Verified: all 10 rows have both \code{NULL}. They are the fossil of a design
where a swap was going to be recorded as a link between two \code{course} rows rather than
as a delete plus an insert. Section~\ref{sec:swap} shows what was built instead, and why
that fossil is a bit of a shame.
\end{honest}

\subsubsection{\code{activity}: the log}

Every action appends a row here. The shipped database has 116 of them. The distinct
\code{activity\_type} values, read out of the file, are exactly:
\code{Add}, \code{Enroll}, \code{Drop}, \code{Swap},
\code{Filter by Day}, \code{Filter by Week}, \code{Export Schedule}, \code{Attendance}.

\begin{honest}{\code{activity.course\_id} is a \code{VARCHAR} that stores prose}
The column is declared \code{Column(String, ForeignKey('course\_list.course\_list\_id'))}.
It is a \emph{text} column claiming to be a foreign key to an \emph{integer} primary key.
It is not used as an identifier at all. These are real, distinct values read out of the
shipped table:

\begin{center}\ttfamily\small
'201' \quad '201, 202' \quad '201, 202, 205, 310, 330' \quad '201 -> 201' \quad 'None' \quad NULL
\end{center}

So the column variously holds one id, a comma-joined list of ids from a multi-course
enroll, an arrow expression from a swap, the four-character string \code{'None'}, and a
genuine SQL \code{NULL}. \code{course\_section} is the same story: it holds \code{'L1'},
\code{'L1, S1'}, \code{'L3 -> L2'}, and for attendance rows an entire sentence.

The log is therefore a \textbf{denormalised, human-readable audit trail}, not a queryable
relation. You cannot ask ``how many times was course 201 dropped'' without parsing strings.
The honest framing: it was built to be rendered into a table widget, and it does that job,
but calling it a foreign key in the model is a false claim the schema makes about itself.

Note the \code{'None'} strings specifically: current code writes
\code{str(course\_id) if course\_id else None}, which yields a real \code{NULL} for a
missing id. So those \code{'None'} rows were written by an \emph{earlier} version of the
code that called \code{str()} unconditionally. The shipped database is a historical
artifact spanning more than one generation of this function. That is a fine thing to
notice out loud; it shows you read the data and not just the code.
\end{honest}

\subsubsection{\code{filter}: declared, created, never used}

\begin{honest}{The \code{filter} table is entirely dead}
\code{Filter} is a full SQLAlchemy model with \code{by\_day} and \code{by\_week} booleans
and a relationship back to \code{Student}. \code{create\_all} duly creates the table in
SQLite. Then: no module imports it, no query touches it, and it holds \textbf{0 rows} in
the shipped database. Verified both ways, by grep and by row count.

What happened is visible in \file{view\_schedule.py}: the day/week toggle state ended up
living in a plain Python attribute, \code{self.filter\_mode}, which is set to
\code{"Day"} at construction and flipped by the buttons. It is never persisted. The
\code{filter} table is what persisting it would have looked like. The screen instead
records each toggle as an \code{activity} row of type \code{Filter by Day}, which is
logging, not state.
\end{honest}

\subsection{The foreign keys do not work, and cannot}
\label{sec:fkbroken}

This is the most substantial finding in the schema, and it was established by running
SQLite against the shipped file, not by reading.

\begin{keyidea}{SQLite does not enforce foreign keys unless you ask it to}
Almost every other database enforces foreign keys by default. SQLite does not: for backward
compatibility, foreign-key enforcement is \textbf{off} unless the connection issues
\code{PRAGMA foreign\_keys = ON}. Verified on the shipped file: the pragma reads
\code{0}. Nothing in \file{student\_database.py} or \file{queries.py} turns it on. So every
\code{ForeignKey} in the model is, at runtime, pure documentation.
\end{keyidea}

Demonstrated directly against a copy of the database, with enforcement at its default:

\begin{lstlisting}[language=SQL,caption={A row referencing a student who does not exist and a course that does not exist. Accepted without complaint.}]
INSERT INTO course VALUES (99999, 88888, 'ZZ', 0, NULL, NULL, NULL);
-- ACCEPTED. No such course_list_id 99999. No such std_id 88888.
\end{lstlisting}

Now the twist. You might think the fix is one line: turn the pragma on. It is not, because
the foreign keys as declared are \emph{structurally invalid} and SQLite will refuse them:

\begin{lstlisting}[language=SQL,caption={With enforcement enabled, the schema itself errors.}]
PRAGMA foreign_keys = ON;
PRAGMA foreign_key_check;
-- OperationalError: foreign key mismatch - "activity" referencing "course_list"

INSERT INTO course VALUES (99999, 88888, 'ZZ', 0, NULL, NULL, NULL);
-- OperationalError: foreign key mismatch - "course" referencing "course"
\end{lstlisting}

\begin{honest}{Why the foreign keys are invalid: they reference half of a composite key}
A foreign key must point at a set of columns that is unique in the parent table. Two of
these do not:

\begin{itemize}
  \item \code{course.course\_id} and \code{activity.course\_id} both reference
        \code{course\_list.course\_list\_id}. But \code{course\_list}'s primary key is the
        \emph{pair} \code{(course\_list\_id, section)}. \code{course\_list\_id} alone is
        \textbf{not unique}: id 201 appears three times, once per section. So the FK
        points at something that does not identify a row, and SQLite calls it a mismatch.
        The correct declaration is a composite FK on \code{(course\_id, section)}.
  \item \code{course.swap\_id} and \code{course.remove\_id} reference
        \code{course.course\_id}, which is likewise only one third of \code{course}'s
        composite PK \code{(course\_id, std\_id, section)}. Same mismatch. The
        commented-out lines just above them in the source show the author already
        wrestling with which table these should point at.
\end{itemize}

The two facts compose into the finding: \textbf{the schema's referential integrity is
broken, and the only reason nothing has ever failed is that SQLite was never asked to
check it.} Turn the pragma on and the application stops working. For a Databases course
project this is the thing to be able to explain, and it is entirely explicable: composite
primary keys make foreign keys harder, SQLAlchemy will happily emit an FK clause it cannot
validate, and SQLite's silence removed the feedback that would have caught it. The fix is
\code{ForeignKeyConstraint(['course\_id','section'], ['course\_list.course\_list\_id','course\_list.section'])}
declared at the table level instead of two column-level \code{ForeignKey} calls.
\end{honest}

\subsection{A relationship that is quietly mis-declared}

\begin{honest}{\code{Course.course\_list} is declared one-to-one and is not}
\file{student\_database.py} declares:

\begin{lstlisting}[language=Python]
course_list = relationship('CourseList', foreign_keys=[course_id],
                           back_populates='course')
\end{lstlisting}

with the source comment ``Define one-to-one relationship between Course and CourseList''.
Because it joins on \code{course\_id} alone and ignores \code{section}, a
\code{Course} row for course 201 matches all \emph{three} \code{course\_list} rows for 201.
This is not theoretical. Running the query layer surfaced it immediately as a real
SQLAlchemy warning:

\begin{lstlisting}
SAWarning: Multiple rows returned with uselist=False for lazily-loaded
attribute 'Course.course_list'
\end{lstlisting}

The application never trips over this because no screen uses the relationship. Every screen
instead calls \code{get\_course\_by\_id\_and\_section(course\_id, section)}, which does the
join by hand with both columns and gets the right answer. So the ORM relationship is dead
weight that is also wrong, and the manual query beside it is correct. That is the same
root cause as Section~\ref{sec:fkbroken}: half of the composite key was dropped.
\end{honest}

\subsection{The import-time side effects}

\begin{lstlisting}[language=Python,caption={The tail of \file{student\_database.py}. All of this runs on \code{import}.}]
engine = create_engine('sqlite:///student.db', echo=True)
...
Base.metadata.create_all(engine)
Session = sessionmaker(bind=engine)
session = Session()
session.close()
\end{lstlisting}

\begin{honest}{Importing any screen creates an engine, writes the schema, and spams SQL}
Four things happen the moment \emph{anything} does \code{import queries}, because
\file{queries.py} imports \file{student\_database}:
\begin{enumerate}
  \item An engine is created against \file{student.db}, bound to the current working
        directory.
  \item \code{echo=True} means SQLAlchemy prints \textbf{every SQL statement it ever
        emits} to the terminal, forever, for the life of the process.
  \item \code{create\_all(engine)} runs, touching the database file and creating any
        missing table.
  \item A module-global \code{session} is opened, then immediately \code{close()}d.
\end{enumerate}
The README already flags this. Two extra things worth knowing that it does not say.

First, the \code{session.close()} on the last line looks like a bug and is not, quite:
SQLAlchemy sessions re-open a connection lazily on next use, so the module-global
\code{session} that every query in \file{queries.py} uses still works. It is confusing and
pointless, but harmless.

Second, and more seriously: \textbf{there is exactly one session for the whole process},
shared by every screen. It is never committed or rolled back on any boundary the user can
see. Combined with the non-atomic swap in Section~\ref{sec:swap}, that shared session is
the thing that makes a half-finished operation stick.
\end{honest}

\subsection{The dead model file}

\begin{honest}{\file{courses\_database.py} defines a second, incompatible \code{Student}}
The README describes \file{courses\_database.py} as ``a seeding helper for the course
catalog''. \textbf{Read the file: it is not.} It contains no course data and no seeding
code at all. What it contains is a second SQLAlchemy \code{Base}, a second engine against
the same \file{student.db}, and a second \code{Student} model whose primary key is called
\code{id} (not \code{std\_id}) and whose \code{dob} is a \code{String} (not a \code{Date}).
Both of those contradict the real \code{student} table. Everything else in the file is
commented out.

Verified: \code{grep} finds no module that imports \code{courses\_database}. It is dead,
its \code{create\_all} is commented out (which is the only reason it has never corrupted
anything), and the README's one-line description of it is wrong. This file should be
deleted. If an interviewer opens it expecting the seeding helper the README promised, you
want to be the one who already knows.
\end{honest}

% =====================================================================
\section{The enrollment rules}
\label{sec:rules}

This is the functional heart of the project. There are four operations the student can
perform on their timetable, and each has its own rules, enforced at different layers and
with different degrees of rigour.

\subsection{The cart: a two-phase commit, done by hand}

\begin{keyidea}{Adding is not enrolling}
The project separates \textbf{choosing} from \textbf{committing}, exactly like an online
shop. \file{add\_courses.py} is the cart: picking a course writes a \code{course} row with
\code{is\_temp=True, type="Add"}. Nothing about the catalogue changes; no seat is taken.
Only when the student goes to \file{enroll.py} and presses Enroll does the flag flip to
\code{is\_temp=False, type=NULL} and the seat count move.

The value of the split is that a cart can be abandoned. The cost is that a
\code{course} row means two different things depending on a boolean, and every query in
\file{queries.py} has to remember to filter on it. Several do; one, as we will see, does
not.
\end{keyidea}

\begin{figure}[H]
\centering
\resizebox{0.94\textwidth}{!}{
\begin{tikzpicture}[node distance=14mm]
  \node[box, fill=white] (none) {\textbf{no row}\\\scriptsize not in cart, not enrolled};
  \node[box, right=52mm of none, fill=accent!15] (cart)
    {\textbf{in cart}\\\code{is\_temp = True}\\\code{type = "Add"}};
  \node[box, right=52mm of cart, fill=okgreen!15] (enr)
    {\textbf{enrolled}\\\code{is\_temp = False}\\\code{type = NULL}};

  \draw[flow] (none) -- node[lbl,above=1pt]{\code{add\_course()}\\\scriptsize Add button,\\\scriptsize add\_courses.py} (cart);
  \draw[flow] (cart) -- node[lbl,above=1pt]{{\scriptsize\ttfamily is\_temp\_false\_}\\{\scriptsize\ttfamily remove\_type\_add("Add")}\\\scriptsize enroll.py, Enroll button} (enr);

  \draw[flow] (cart) to[out=140,in=40,looseness=1.1]
     node[lbl,above]{\code{remove\_course()} \scriptsize Remove button} (none);
  \draw[flow] (enr) to[out=-140,in=-40,looseness=0.55]
     node[lbl,below,pos=0.5]{\code{remove\_course()} \scriptsize Drop button, drop\_courses.py\\
       \scriptsize and the \emph{first half} of a swap} (none);

  \draw[flow, draw=warnred] (enr) to[out=40,in=140,looseness=4]
     node[lbl,above]{\scriptsize swap: enr $\rightarrow$ none $\rightarrow$ cart $\rightarrow$ enr,\\
       \scriptsize in three separate commits (\S\ref{sec:swap})} (enr);

  \node[below=22mm of cart, font=\scriptsize\itshape, text=black!60, align=center]
    {The only two persisted states of an enrollment. Everything else the UI shows\\
     (the catalogue's Open/Closed, the day/week toggle) lives elsewhere.};
\end{tikzpicture}}
\caption{The lifecycle of a \code{course} row. The red self-loop is the swap, and it is the
one transition that is not safe.}
\label{fig:cartsm}
\end{figure}

\subsection{Rule 1: can this course be added?}

The gate is \code{AddCoursesWindow.can\_the\_course\_be\_added(course\_id, start\_time,
end\_time, days)}, backed by a second check inside \code{AllQueries.add\_course}.

\begin{figure}[H]
\centering
\resizebox{0.98\textwidth}{!}{
\begin{tikzpicture}[node distance=9mm]
  \node[box] (start) {Add button clicked\\\scriptsize on the selected catalogue row};
  \node[dec, below=of start, text width=24mm] (incart) {already a row\\for this\\course id\\in the cart?};
  \node[dec, below=11mm of incart, text width=24mm] (enrolled) {already\\\emph{enrolled} in\\this course id?};
  \node[dec, below=11mm of enrolled, text width=24mm] (clash) {timings clash\\with anything\\already held?};
  \node[dec, below=11mm of clash, text width=24mm] (closed) {status\\column reads\\"Closed"?};
  \node[dec, below=11mm of closed, text width=26mm] (dup) {\code{check\_if\_course}\\\code{\_already\_exist}\\(id + section)?};
  \node[box, below=11mm of dup, fill=okgreen!15] (ok) {\textbf{INSERT} course row\\\code{is\_temp=True, type="Add"}\\then log an \code{"Add"} activity};

  \node[box, right=38mm of incart, fill=warnred!12, draw=warnred]
       (r1) {\scriptsize "Cannot add the same course twice"};
  \node[box, right=38mm of enrolled, fill=warnred!12, draw=warnred]
       (r2) {\scriptsize "Cannot add the same course twice"};
  \node[box, right=38mm of clash, fill=warnred!12, draw=warnred]
       (r3) {\scriptsize "Timings clash"};
  \node[box, right=38mm of closed, fill=warnred!12, draw=warnred]
       (r4) {\scriptsize "Course is closed"};
  \node[box, right=38mm of dup, fill=warnred!12, draw=warnred]
       (r5) {\scriptsize "Course could not be added"};

  \draw[flow] (start) -- (incart);
  \draw[flow] (incart) -- node[lbl,left]{no} (enrolled);
  \draw[flow] (enrolled) -- node[lbl,left]{no} (clash);
  \draw[flow] (clash) -- node[lbl,left]{no} (closed);
  \draw[flow] (closed) -- node[lbl,left]{no} (dup);
  \draw[flow] (dup) -- node[lbl,left]{no} (ok);

  \draw[flow] (incart) -- node[lbl,above]{yes} (r1);
  \draw[flow] (enrolled) -- node[lbl,above]{yes} (r2);
  \draw[flow] (clash) -- node[lbl,above]{yes} (r3);
  \draw[flow] (closed) -- node[lbl,above]{yes} (r4);
  \draw[flow] (dup) -- node[lbl,above]{yes} (r5);

  \node[extbox, left=14mm of clash, text width=38mm]
    {\scriptsize checks against \emph{all} the student's rows, cart included:\\[2pt]
     {\tiny\ttfamily get\_days\_timings\_of\_}\\{\tiny\ttfamily all\_courses\_by\_id\_}\\{\tiny\ttfamily temp\_false()}};
  \draw[dflow] (clash) -- (clash-|-3,0);
\end{tikzpicture}}
\caption{The add gate, in the order the code actually evaluates it. Note the first four
tests live in the window class and only the last lives in the query layer.}
\label{fig:addflow}
\end{figure}

Three things about that flowchart are worth stopping on.

\begin{honest}{The ``already added'' rule is by course id, so you cannot take two sections
of one course, and the duplicate check disagrees with itself}
The two top gates compare \textbf{course id only}, ignoring section. So having 201/L1 in
your cart blocks you from adding 201/L2. That is almost certainly the intent (you take one
section of a course), and the message ``Cannot add the same course twice'' says so.

But the last gate, \code{check\_if\_course\_already\_exist(course\_id, section)} inside
\code{AllQueries.add\_course}, compares \textbf{id and section together}. So the query layer
would happily allow 201/L1 and 201/L2 to coexist; only the UI stops it. And the database
agrees with the query layer, because \code{course}'s primary key includes \code{section}.

So the rule ``one section per course'' is enforced in exactly one place, the weakest one:
a window class. Nothing in the schema or the data layer encodes it. Anything that inserts
by another path, including \code{swap\_courses}, bypasses it entirely.
\end{honest}

\begin{honest}{\code{get\_days\_timings\_of\_all\_courses\_by\_id\_temp\_false} does not
filter on \code{is\_temp}}
Read the name, then read the body:
\begin{lstlisting}[language=Python]
def get_days_timings_of_all_courses_by_id_temp_false(self):
    std_id = self.glob_id
    course_list_query = session.query(Course).filter_by(std_id=std_id).all()
\end{lstlisting}
No \code{is\_temp=False}. Its sibling one method above, \code{get\_timings\_of\_all\_
courses\_by\_id}, does filter. So the method whose name promises to exclude cart rows is
the one that includes them.

Here is the subtle part, and it is worth getting right rather than just calling it a bug:
\textbf{the unfiltered behaviour is the behaviour the cart actually needs}. When you are
adding a third course, you want it checked against your enrolled courses \emph{and} the two
already in your cart. If it filtered as its name promises, you could fill a cart with three
mutually clashing courses and enroll in all of them. So the code is right and the name is
a lie. That is a much better thing to say in an interview than ``there is a bug''. The fix
is a rename, to something like \code{get\_days\_timings\_of\_all\_courses\_by\_id}, and
a comment saying why the cart counts.
\end{honest}

\begin{honest}{The clash check runs before the closed check, and both run before either matters}
In \code{add\_course}, \code{can\_the\_course\_be\_added(...)} is called on line 216 and
its result is stored, then the \code{status == "Closed"} test runs on line 219 and returns
first. So on a closed course that also clashes, the user sees the ``Timings clash'' warning
popup (raised from inside the clash check) and \emph{then} the ``Course is closed'' popup.
Two dialogs for one click. Cosmetic, but it is the kind of thing a demo surfaces.
\end{honest}

\subsection{Rule 2: the clash detector}
\label{sec:clash}

This is the most algorithmically interesting code in the project. It answers: does the
candidate section overlap in time, on a shared day, with anything the student already
holds?

\begin{jargon}{Interval overlap}
Two time intervals overlap if one starts before the other ends and ends after the other
starts. The naive way to test it is to enumerate cases: does A's start fall inside B? does
A's end fall inside B? does A completely contain B? Those three cases are exhaustive, and
they are exactly the three conditions this code tests.
\end{jargon}

The code works in two steps: reduce times to a comparable number, then compare.

\begin{lstlisting}[language=Python,caption={\file{queries.py}: times become integers.}]
def time_to_seconds(self, t):
    return t.hour * 3600 + t.minute * 60 + t.second

def get_days_list(self, days):
    return [days[x:x + 2] for x in range(0, len(days), 2)]
\end{lstlisting}

\code{get\_days\_list} is how \code{"MoWe"} becomes \code{['Mo', 'We']}: it slices the
string into two-character chunks. This is the entire day-encoding scheme, and it is why
the catalogue must never contain a day code of any other length.

\begin{figure}[H]
\centering
\resizebox{0.95\textwidth}{!}{
\begin{tikzpicture}[xscale=1.25]
  % axis
  \draw[->, thick, draw=inkblue] (0,0) -- (11.4,0) node[right,font=\scriptsize] {time (seconds since midnight)};
  \foreach \x/\l in {0/08:30, 2/10:00, 4/11:30, 6/13:00, 8/14:30, 10/16:00}
    \draw[draw=linegrey] (\x,-0.12) -- (\x,0.12) node[below=3pt,font=\scriptsize\ttfamily,text=black!60] {\l};

  % existing course
  \draw[fill=inkblue!18, draw=inkblue] (4,0.45) rectangle (5.7,1.05);
  \node[font=\scriptsize, text=inkblue] at (4.85,0.75) {already held};
  \node[font=\tiny\ttfamily, text=black!60] at (4.85,1.3) {start\_timings[i] .. end\_timings[i]};

  % case A
  \draw[fill=warnred!18, draw=warnred] (5.0,-1.15) rectangle (6.8,-0.55);
  \node[font=\scriptsize\bfseries, text=warnred] at (5.9,-0.85) {A};
  \node[font=\scriptsize, text=warnred, anchor=west] at (7.0,-0.85) {new start inside:};
  \node[font=\tiny\ttfamily, text=black!60, anchor=west] at (9.3,-0.85)
    {start[i] <= new\_start <= end[i]};

  % case B
  \draw[fill=warnred!18, draw=warnred] (2.6,-2.05) rectangle (4.4,-1.45);
  \node[font=\scriptsize\bfseries, text=warnred] at (3.5,-1.75) {B};
  \node[font=\scriptsize, text=warnred, anchor=west] at (7.0,-1.75) {new end inside:};
  \node[font=\tiny\ttfamily, text=black!60, anchor=west] at (9.3,-1.75)
    {start[i] <= new\_end <= end[i]};

  % case C
  \draw[fill=warnred!18, draw=warnred] (3.2,-2.95) rectangle (6.5,-2.35);
  \node[font=\scriptsize\bfseries, text=warnred] at (4.85,-2.65) {C};
  \node[font=\scriptsize, text=warnred, anchor=west] at (7.0,-2.65) {new envelops held:};
  \node[font=\tiny\ttfamily, text=black!60, anchor=west] at (9.3,-2.65)
    {new\_start <= start[i] and new\_end >= end[i]};

  % case D no clash
  \draw[fill=okgreen!18, draw=okgreen] (8.4,-3.85) rectangle (10.2,-3.25);
  \node[font=\scriptsize\bfseries, text=okgreen] at (9.3,-3.55) {D};
  \node[font=\scriptsize, text=okgreen, anchor=west] at (10.4,-3.55) {no overlap};
  \node[font=\tiny\itshape, text=black!60, anchor=west] at (0,-4.3)
    {(and: only tested at all if the two courses share a two-character day code)};

  \draw[dashed, draw=inkblue!40] (4,1.05) -- (4,-4.0);
  \draw[dashed, draw=inkblue!40] (5.7,1.05) -- (5.7,-4.0);
\end{tikzpicture}}
\caption{The three overlap cases the clash detector tests, drawn against one already-held
section. A, B and C together are exhaustive: any overlap of two intervals falls into at
least one of them.}
\label{fig:clash}
\end{figure}

\begin{pseudocode}{The clash rule, as implemented in \code{can\_the\_course\_be\_added}}
\Step{new\_start, new\_end := time\_to\_seconds(candidate start, end)}
\Step{for each course the student already holds (enrolled OR in cart), as i:}
\Indent{shared := any 2-char day code of the candidate}
\Indent{\phantom{shared := }appears in the day codes of course i}
\Indent{if shared:}
\Indent{\hspace*{2em}if start[i] <= new\_start <= end[i]      -> CLASH}
\Indent{\hspace*{2em}if start[i] <= new\_end   <= end[i]      -> CLASH}
\Indent{\hspace*{2em}if new\_start <= start[i] and new\_end >= end[i] -> CLASH}
\Step{return no clash}
\end{pseudocode}

\begin{honest}{Two real defects in the clash rule}
\textbf{First, the day test is inside the wrong loop and re-derived every iteration.} The
code reads:
\begin{lstlisting}[language=Python]
for i in range(len(start_timings)):
    days_list = self.all_queries.get_days_list(days)              # candidate
    previous_days_list = self.all_queries.get_days_list(all_days[i])  # held
    for day in range(len(days_list)):
        if days_list[day] in previous_days_list:
            ... three overlap tests ...
\end{lstlisting}
\code{days\_list} does not depend on \code{i}, so it is recomputed on every iteration of
the outer loop. Harmless, wasteful, and a sign the loop grew rather than being designed.
More notably, the three overlap tests sit inside the \emph{inner} day loop, so for a
candidate on \code{MoWe} clashing with a held course on \code{MoWe}, they run twice and
return on the first. Correct, but only by accident of the early return.

\textbf{Second, and this one is a genuine bug: touching intervals are treated as a clash.}
All three tests use \code{<=}. A candidate starting at exactly 12:45 against a held course
ending at exactly 12:45 satisfies \code{start[i] <= new\_start <= end[i]} and is rejected.
Look at the real catalogue: 201/L1 runs 11:30 to 12:45 on MoWe, and back-to-back classes
are the normal case in a university timetable. This rule forbids them. A strict
\code{<} on the boundary comparisons is the fix. This is easy to demonstrate live and
worth knowing before someone else demonstrates it to you.
\end{honest}

\begin{honest}{The clash rule is only ever consulted by the UI}
\code{AllQueries.add\_course} does not call the clash checker. It is called only from
\code{AddCoursesWindow.can\_the\_course\_be\_added} and
\code{SwapCoursesWindow.can\_the\_course\_be\_added}, which are two separate,
near-identical copies of the same 40 lines in two different files. So: a clash-free
timetable is a property the \emph{windows} maintain, not a property the data layer or the
schema guarantees, and the rule is duplicated so it can drift. It has already begun to:
the two copies call \emph{different} query methods (the add path uses the unfiltered one
that includes the cart; the swap path uses the filtered one with the exclusion arguments).
\end{honest}

\subsection{Rule 3: capacity, and what \code{current\_capacity} actually means}
\label{sec:capacity}

\begin{keyidea}{\code{current\_capacity} counts \emph{empty seats}, not filled ones}
This is the single most counter-intuitive thing in the codebase, and the column name
actively works against you. Read \code{update\_current\_capacity}:
enrolling calls it with \code{"Subtract"}; dropping calls it with \code{"Add"}. So the
number goes \emph{down} as students join. It is \textbf{remaining seats}.

The shipped data confirms it: untouched sections read \code{current\_capacity = 50}, equal
to \code{max\_capacity}, and the sections with enrollments read 47. A name like
\code{seats\_remaining} would have removed all doubt.
\end{keyidea}

\begin{lstlisting}[language=Python,caption={\file{queries.py}, the whole capacity rule.}]
def update_current_capacity(self, course_id, section, add_or_subtract):
    course = session.query(CourseList).filter_by(course_list_id=course_id,
                                                 section=section).first()
    if add_or_subtract == "Add":
        course.current_capacity += 1
    elif add_or_subtract == "Subtract":
        course.current_capacity -= 1

    if course.current_capacity == 0:
        course.status = False
    elif course.current_capacity > 0:
        course.status = True
    else:
        print("Error: current_capacity is negative")
        return False
    ...
\end{lstlisting}

So \code{status} is a derived flag: false (``Closed'') exactly when no seats remain. The
catalogue screens hide any section whose \code{status} is false, which is how a full
section disappears from the list.

This was verified by running the real method against a copy of the database:

\begin{lstlisting}[caption={Executed against a copy of \file{student.db}. Output is real.}]
201/L2 before:                 max=50, current=50, status=True
after an ENROLL  (Subtract):   max=50, current=49, status=True
after two DROPs  (Add, Add):   max=50, current=51, status=True   <-- exceeds max
\end{lstlisting}

\begin{honest}{\code{max\_capacity} is never read by a single line of code}
Verified by grep across every \file{.py} and \file{.ui}: the only occurrence of
\code{max\_capacity} anywhere in the project is the column declaration itself. It is
written into the seed data and never consulted again. The ceiling on a section is not
enforced; only the floor is (\code{current\_capacity == 0} closes it).

The consequence is not hypothetical. \textbf{The shipped database already contains two
sections that are over capacity}: 330/L3 and 310/L2 both have
\code{current\_capacity = 51} against \code{max\_capacity = 50}. Somebody dropped a course
they had never enrolled in (or the seat count was decremented once and incremented twice),
and nothing stopped the counter climbing past the maximum. Those rows are in the file you
would demo from.

The fix is two lines in the same function that already has the ``negative'' branch: refuse
to increment past \code{max\_capacity}, and refuse to decrement below zero \emph{before}
committing rather than after. Note the existing negative branch returns \code{False} only
\emph{after} having already mutated the object, and its caller in \file{enroll.py} shows
``Invalid capacity'' and returns, leaving the session holding a dirty object. It happens
not to be committed on that path, but nothing rolls it back either.
\end{honest}

\begin{honest}{Capacity is not decremented where you would expect}
Adding to the cart deliberately does \emph{not} take a seat: the
\code{update\_current\_capacity} call in \file{add\_courses.py} is present but commented
out. Only Enroll takes the seat. This is the right call (a cart should not hold
inventory), but it also means there is \textbf{no reservation}: two students can both pass
the ``Open'' check, both fill their carts with the last seat, and both enroll. The second
one drives \code{current\_capacity} to $-1$, at which point the negative branch fires,
\code{status} is left stale, and the user sees ``Invalid capacity''. For a single-user
coursework demo this never happens. It is the honest answer to ``what breaks if two
students use this at once'', and the answer is: the seat count, immediately, because there
is no locking and no constraint.
\end{honest}

\subsection{Rule 4: the swap}
\label{sec:swap}

A swap moves the student from one section of a course to another section of the
\emph{same} course. The UI enforces exactly that:

\begin{lstlisting}[language=Python,caption={\file{swap\_courses.py}. Note the comment describes the opposite of the code.}]
# if section and course_id same then show messagebox and return
if course_id != course_id_enrolled:
    QMessageBox.warning(self, "Error", "Cannot swap different courses.")
    return

if course_id == course_id_enrolled and section == section_enrolled:
    QMessageBox.warning(self, "Error", "Cannot swap the same course.")
    return
\end{lstlisting}

So the two rules are: the courses must match, and the sections must differ. The Select
button already filters the second table to sections of the chosen course, via
\code{get\_all\_courses\_by\_std\_id\_}\-\code{and\_course\_id}, so the first rule mostly
guards against the user clicking around it.

\begin{figure}[H]
\centering
\resizebox{\textwidth}{!}{
\begin{tikzpicture}[node distance=6mm]
  % lifelines
  \node[box, fill=accent!15] (ui) at (0,0) {SwapCoursesWindow\\\scriptsize swap\_courses.py};
  \node[box, fill=softgrey] (q) at (5.4,0) {AllQueries.swap\_courses\\\scriptsize queries.py};
  \node[store] (db) at (11.0,0) {student.db};

  \draw[draw=linegrey, thick] (ui) -- (0,-9.4);
  \draw[draw=linegrey, thick] (q) -- (5.4,-9.4);
  \draw[draw=linegrey, thick] (db) -- (11.0,-9.4);

  \draw[flow] (0,-1.0) -- node[lbl,above]{swap\_courses(new, section, old, old\_section)} (5.4,-1.0);

  \draw[flow] (5.4,-2.0) -- node[lbl,above]{\textbf{1. remove\_course(old)} -> DELETE} (11.0,-2.0);
  \draw[flow] (11.0,-2.5) -- node[lbl,above]{\textbf{COMMIT} \scriptsize the old course is now gone, on disk} (5.4,-2.5);

  \draw[flow] (5.4,-3.5) -- node[lbl,above]{\textbf{2. add\_course(new)} -> INSERT is\_temp=True} (11.0,-3.5);

  \draw[flow, draw=warnred] (5.4,-4.6) -- node[lbl,above,text=warnred]{\textbf{if step 2 fails:} return False / None,\\
       \scriptsize nothing is rolled back} (0,-4.6);

  \draw[flow] (5.4,-5.6) -- node[lbl,above]{\textbf{3. is\_temp\_false\_single\_course(new)} -> UPDATE} (11.0,-5.6);
  \draw[flow] (11.0,-6.1) -- node[lbl,above]{\textbf{COMMIT} \scriptsize now enrolled} (5.4,-6.1);
  \draw[flow] (5.4,-6.6) -- node[lbl,above]{True} (0,-6.6);

  \draw[flow] (0,-7.6) -- node[lbl,above]{4. update\_current\_capacity(old, "Add") \scriptsize give the seat back} (11.0,-7.6);
  \draw[flow] (0,-8.4) -- node[lbl,above]{5. update\_current\_capacity(new, "Subtract") \scriptsize take the new seat} (11.0,-8.4);
  \draw[flow] (0,-9.2) -- node[lbl,above]{6. log a "Swap" activity: "201 -> 201", "L3 -> L2"} (11.0,-9.2);

  \node[draw=warnred, fill=warnred!8, rounded corners=2pt, text width=52mm,
        font=\scriptsize, align=left] at (16.5,-4.0)
    {\textbf{The unsafe window.}\\[2pt]
     Step 1 commits before step 2 runs. There is no transaction spanning them. If step 2
     fails, the student has \emph{lost} the old course and gained nothing, the seat counts
     in steps 4 and 5 never run, and the UI reports\\
     \textit{``Course could not be swapped.''}\\[2pt]
     Which is false. Something very much happened.};

\end{tikzpicture}}
\caption{The swap, step by step, showing where the atomicity is missing. Verified by
execution, below.}
\label{fig:swapseq}
\end{figure}

\begin{jargon}{Atomicity and transactions}
An operation is \textbf{atomic} if it either happens completely or not at all, with no
observable in-between. Databases provide this with a \textbf{transaction}: you begin one,
make several changes, and \code{COMMIT} them as a single indivisible act, or
\code{ROLLBACK} to undo the lot. A swap is the textbook case: ``delete the old, insert the
new'' must be atomic, or a crash between the two loses the student's course. This is the
`A' in ACID, and it is a core Databases-course concept, which makes it the most important
finding in this document.
\end{jargon}

\begin{honest}{The swap is not atomic, and it destroys data. Verified by running it.}
\code{AllQueries.swap\_courses} calls \code{remove\_course()} first, which does
\code{session.delete(course)} followed by its own \code{session.commit()}. That commit is
unconditional and it is outside the \code{try} block that guards the add. Only then does
it attempt \code{add\_course()}.

This was tested directly against a copy of the shipped database, by asking for a swap onto
a section that does not exist so that step 2 must fail:

\begin{lstlisting}[caption={Executed against a copy of \file{student.db}. Output is real.}]
enrolled rows for student 20180 before:
    [(201,'L1'), (202,'S2'), (205,'L3'), (330,'L3'), (310,'L2')]

swap_courses(new=999/'ZZ' (nonexistent), old=201/'L1')  ->  returned: None

enrolled rows for student 20180 after:
    [(202,'S2'), (205,'L3'), (330,'L3'), (310,'L2')]

LOST:   [(201, 'L1')]
GAINED: []
\end{lstlisting}

The student's enrollment in 201/L1 was deleted and committed, the replacement was never
created, and the function returned \code{None}. \code{None} is falsy, so the caller in
\file{swap\_courses.py} takes the \code{if not can\_be\_added:} branch and shows the user
\textit{``Course could not be swapped.''} The user is told nothing happened. They have in
fact just lost a course, permanently, and the seat is not even returned to the pool because
steps 4 and 5 are below the early return.

Note also the bare \code{except:} inside \code{swap\_courses} ends in \code{return} with no
value, which is why the function returns \code{None} rather than \code{False}. Two
different falsy failure signals from one function.

\textbf{The fix is small and you should be able to state it instantly}: wrap steps 1 to 3
in one transaction. Do not commit inside \code{remove\_course} when it is called from a
swap; do the delete, do the insert, and commit once, with \code{session.rollback()} in the
except. SQLAlchemy gives you this for free with \code{with session.begin():}. The reason
it was not done is visible in the code: \code{remove\_course} and \code{add\_course} were
each written to be self-contained, committing operations for their own screens, and
\code{swap\_courses} reuses them as-is. Composing two atomic operations does not give you
an atomic operation. That sentence is the whole lesson.

And this is what the dead \code{swap\_id} / \code{remove\_id} columns were for. The
original design was going to \emph{link} the new row to the old one rather than delete and
re-insert. That design would not have had this bug.
\end{honest}

\subsection{Rule 5: the drop}

The simplest path, in \file{drop\_courses.py}: confirm with a Yes/No dialog, delete the
\code{course} row, return the seat with \code{update\_current\_capacity(..., "Add")},
remove the table row, log a \code{"Drop"} activity. There is no rule about
\emph{whether} you may drop; you always may.

\begin{honest}{The drop screen's column headers do not line up with its data}
\code{setup\_drop\_courses\_table} sets nine headers:
\code{["Course ID", "Course Name", "Section", "Instructor", "Days", "Start Time", "End Time", "Room", "Status"]}.
Then it fills only \emph{eight} columns, because it merges start and end into one:

\begin{center}
\begin{tabular}{@{}clll@{}}
\toprule
col & header says & the code actually puts there & \\
\midrule
4 & Days & days & ok \\
5 & Start Time & \code{"11:30:00 - 12:45:00"} & both times in one cell \\
6 & End Time & \textbf{room} & wrong \\
7 & Room & \textbf{status} & wrong \\
8 & Status & \textbf{nothing} & empty column \\
\bottomrule
\end{tabular}
\end{center}

So on the drop screen the room number appears under a heading that says ``End Time'', the
word ``Open'' appears under ``Room'', and the last column is blank. This is a visible
defect on a screen a demo will certainly open. The functional part is unaffected, since
\code{drop\_course} reads columns 0 and 2, which are still correct.
\end{honest}

\subsection{The enroll screen, and a bug that cancels itself out}
\label{sec:checkstate}

\file{enroll.py} is the checkout. It lists everything in the cart
(\code{get\_all\_courses\_by\_id\_and\_type("Add")}), refuses if any row reads ``Closed'',
then flips the whole cart to enrolled, then decrements a seat for each row.

\begin{lstlisting}[language=Python,caption={\file{enroll.py}, \code{enroll\_action}. The order matters.}]
for row in range(self.schedule_table_widget.rowCount()):
    if self.schedule_table_widget.item(row, 8).text() == "Closed":
        QMessageBox.about(self, "Enroll", "You cannot enroll in a closed course")
        return

self.all_queries.is_temp_false_remove_type_add("Add")     # <-- commits EVERYTHING
QMessageBox.about(self, "Enroll", "You have successfully enrolled ...")

course_ids = []
for row in range(self.schedule_table_widget.rowCount()):
    if self.schedule_table_widget.item(row, 0).checkState():   # <-- always True
        course_ids.append(self.schedule_table_widget.item(row, 0).text())
\end{lstlisting}

\begin{honest}{\code{if item.checkState():} is always true in PyQt6. Verified.}
The intent is clearly ``only the rows the user ticked''. It does not work, and the reason
is a genuine PyQt5-to-PyQt6 migration trap.

In PyQt5, \code{checkState()} returned a plain integer, and \code{Qt.Unchecked} was
\code{0}, so \code{if item.checkState():} was false for an unticked row. PyQt6 replaced
those integer constants with real Python enums, and \textbf{a Python enum member is truthy
regardless of its value} unless the enum defines otherwise. \code{Qt.CheckState} is a plain
\code{enum.Enum}, not an \code{IntEnum}.

Verified by installing PyQt6 6.11 in a throwaway environment and asking:
\begin{lstlisting}
>>> repr(Qt.CheckState.Unchecked)
<CheckState.Unchecked: 0>
>>> bool(Qt.CheckState.Unchecked)
True
\end{lstlisting}
So the condition is true for every row, ticked or not. The whole cart is always collected.

Now the interesting part: \textbf{this bug has no visible effect}, for two compounding
reasons. First, \code{is\_temp\_false\_remove\_type\_add("Add")} on the line above has
already enrolled the entire cart in the database before the loop runs; the loop only
decides what gets logged and whose seat gets taken. Second, no code ever calls
\code{setCheckState} on these items and the table is set to \code{NoEditTriggers}, so no
checkbox is ever shown and the user could not tick anything anyway. The intended feature
does not exist in the UI, and the code that would have implemented it is broken in a way
that happens to produce the behaviour the rest of the screen assumes: enroll everything.

Two wrongs make a right here. It is a much better story than a bug: it shows you can trace
a defect to the exact language-level semantics that caused it, and then judge honestly that
it does not currently matter. What it \emph{does} mean is that the dead code should be
deleted, because the day someone adds real checkboxes, this silently enrolls the unticked
ones too.
\end{honest}

\begin{honest}{The success dialog fires before the work that can fail}
``You have successfully enrolled in the selected courses'' is shown immediately after the
\code{is\_temp} flip and \emph{before} the capacity loop. If
\code{update\_current\_capacity} then returns false, the user gets ``Invalid capacity''
straight after being congratulated. This is the same pattern the attendance module
deliberately avoids (Section~\ref{sec:attendance}), which makes the inconsistency worth
noting: the newer code learned the lesson, the older code did not.
\end{honest}

% =====================================================================
\section{The data-access layer}

\code{AllQueries} in \file{queries.py} is 410 lines and about 25 methods, and it is the
project's best structural idea: no window ever writes a query, so the screens agree on a
contract. The README says this was deliberate, to keep a two-person team's merges
tractable, and the code bears that out.

\subsection{\code{glob\_id}: how the layer knows who you are}

Every method that touches student-specific data begins the same way:

\begin{lstlisting}[language=Python]
def get_all_courses_by_id(self):
    std_id = self.glob_id
    ...
\end{lstlisting}

\code{glob\_id} is set in \code{\_\_init\_\_} to a hardcoded \code{20180} and then
overwritten by every screen with
\code{self.all\_queries.update\_std\_id(self.all\_queries.get\_std\_id\_by\_email(self.email))}.

\begin{honest}{The README overstates this one, and the truth is more interesting}
The README's ``What I learned'' says: \textit{``\code{AllQueries.glob\_id} (defaulting to a
hardcoded 20180) is set by whichever window ran last, so the query layer's correctness
depends on call order.''}

Read the code: \code{glob\_id} is assigned in \code{\_\_init\_\_}, so it is an
\textbf{instance} attribute, not a class attribute. Every screen constructs its \emph{own}
\code{AllQueries()} and immediately sets its own \code{glob\_id} from its own
\code{self.email}. So screens do not clobber each other's student id, and the stated
call-order hazard does not exist as described. Correct the README rather than repeat this
claim in an interview: the fix is one word, ``instance'', and being the person who caught
it is better than being the person who recited it.

The real hazards are different and both worth defending:

\textbf{1. The default 20180 is a live student.} It is Alice, a real row in the shipped
database. Any \code{AllQueries()} constructed without a subsequent \code{update\_std\_id}
silently operates \emph{as Alice}. Verified: constructing \code{AllQueries()} and reading
\code{glob\_id} gives \code{20180} with no login of any kind. Every screen's
\code{\_\_init\_\_} also defaults \code{email} to \code{alice20180@st.habib.edu.pk} when
none is passed, which is what the \code{if \_\_name\_\_ == "\_\_main\_\_"} block at the
bottom of every module does. So each screen is independently runnable as Alice, with no
password and no verification code. That is the actual security-shaped hole, and it is a
debugging convenience that shipped.

\textbf{2. The \emph{session} genuinely is shared.} \code{session} is a module-global in
\file{student\_database.py}, created once at import. Every \code{AllQueries} instance uses
the same one. So the thing the README worried about is real, it is just one level down: not
the id, the session. That is what makes the swap's partial commit in
Section~\ref{sec:swap} durable.
\end{honest}

\subsection{Error handling, uniformly}

Nearly every mutating method follows one shape:

\begin{lstlisting}[language=Python]
session.add(new_course)
try:
    session.commit()
except:
    print("Error: " + str(sys.exc_info()[0]))
    print(traceback.format_exc())
    session.rollback()
    return False
return True
\end{lstlisting}

\begin{honest}{Bare \code{except:} everywhere, and \code{print} as the logging strategy}
\code{except:} with no exception class catches \emph{everything}, including
\code{KeyboardInterrupt} and \code{SystemExit}. It should be \code{except Exception:} at
minimum, and really \code{except SQLAlchemyError:}. The failure is reported by printing to
a terminal that a GUI user is not looking at, and returning \code{False}, which some
callers check and some ignore.

The consistency is genuinely a plus: one shape, applied everywhere, is easier to fix than
twenty bespoke ones. Replacing \code{print} with the \code{logging} module and the bare
\code{except} with a typed one is a mechanical change across about a dozen methods. The
README already lists this as a follow-up.
\end{honest}

\begin{honest}{\code{remove\_course} does not check that the row exists}
\begin{lstlisting}[language=Python]
course = session.query(Course).filter_by(...).first()
session.delete(course)      # <-- course may be None
try:
    session.commit()
\end{lstlisting}
The \code{delete} is \emph{outside} the try. If the row is not found, \code{course} is
\code{None} and \code{session.delete(None)} raises
\code{UnmappedInstanceError}, which propagates past the try that was supposed to guard it,
out of \code{remove\_course}, out of \code{swap\_courses}, and into the window's own bare
\code{except}, where it becomes a printed traceback and a silent return. Two of the three
callers happen to pass ids read out of a table populated from the database, so the row
does exist. It is guarded by luck.
\end{honest}

\subsection{Two methods that do the same job}

Two methods insert an \code{Activity} row:

\begin{center}
\code{add\_data\_to\_activity\_log(std\_id, activity\_type, course\_id, section)}\\[2pt]
\code{insert\_activity(std\_id, activity\_type, description, date\_val, time\_val)}
\end{center}

\noindent There is also a third, \code{insert\_attendance\_activity}, which formats a
string and inserts an \code{Activity} row.

\begin{honest}{Three activity-insert methods; one is dead}
\code{insert\_attendance\_activity} is never called by anything. Verified by grep:
\file{attendance.py} calls \code{insert\_activity}. The section comment above it in
\file{queries.py} announces
\code{\# -------------- NEW METHOD FOR ATTENDANCE --------------}, so it is the first
attempt, superseded by \code{insert\_activity} directly below it, and never removed. Its
docstring even contains a smart apostrophe, marking it as pasted from somewhere else.

The difference between the two survivors is that \code{add\_data\_to\_activity\_log} stamps
the time itself with \code{datetime.now()}, while \code{insert\_activity} takes date and
time as \emph{strings} and parses them with \code{strptime}, falling back to
\code{now()} if parsing fails. The caller in \file{attendance.py} formats \code{now()}
into strings, passes them, and has them parsed back. That is a round trip through text for
no reason. One method taking a \code{datetime} would have replaced all three.
\end{honest}

Also note \code{add\_data\_to\_activity\_log} contains a query whose result is discarded:

\begin{lstlisting}[language=Python]
course = None
if course_id is not None:
    course = session.query(CourseList).filter_by(course_list_id=course_id).first()
# `course` is never used again.
\end{lstlisting}

A wasted \code{SELECT} on every logged action.

\subsection{The destructive \code{\_\_main\_\_} block}

\begin{honest}{Running \code{python queries.py} wipes your data}
\begin{lstlisting}[language=Python]
if __name__ == "__main__":
    all_queries = AllQueries()
    all_queries.delete_all_data_activity_log()
    all_queries.delete_everything_in_course_table()
    print(all_queries.get_current_day())
\end{lstlisting}
That deletes \textbf{every row} of \code{activity} and \textbf{every row} of \code{course},
for every student, with no confirmation, as the price of running the file. It is obviously
a development reset script that was left in the shipped module. It should be a separate
\file{reset\_db.py}, or at minimum behind an \code{argv} check. Note also that
\code{delete\_all\_data\_activity\_log} and \code{delete\_everything\_in\_course\_table}
ignore \code{glob\_id} entirely: they are not scoped to a student despite living on an
object whose entire purpose is to scope queries to a student.
\end{honest}

% =====================================================================
\section{Login and the verification flow}

This is one of the areas the README records as Salman's, and it is the most carefully
written code in the project.

\begin{figure}[H]
\centering
\resizebox{0.96\textwidth}{!}{
\begin{tikzpicture}[node distance=9mm]
  \node[box] (a) {user types email + password};
  \node[dec, below=of a, text width=26mm] (b) {\code{get\_password\_}\\\code{by\_email}\\returns None?};
  \node[dec, below=12mm of b, text width=26mm] (c) {\code{password ==}\\\code{user\_password}?};
  \node[box, below=12mm of c, fill=accent!12] (d) {\textbf{VerificationForm} opens\\
     \scriptsize \code{random.randint(100000, 999999)}};
  \node[dec, below=of d, text width=30mm] (e) {\code{ATTENDANCE\_SENDER\_}\\\code{EMAIL} and
     \code{\_PASSWORD}\\both set in the env?};
  \node[box, below left=12mm and 6mm of e, fill=softgrey] (f)
    {\textbf{fallback}\\\scriptsize print the code to the terminal\\
     \scriptsize + warn "Email Not Configured"\\\scriptsize \textbf{login still completes}};
  \node[box, below right=12mm and 6mm of e, fill=softgrey] (g)
    {\textbf{send}\\\scriptsize smtplib -> smtp.gmail.com:587\\\scriptsize starttls, login, sendmail};
  \node[dec, below=26mm of e, text width=26mm] (h) {typed code ==\\generated code?};
  \node[box, below=12mm of h, fill=okgreen!15] (i) {\textbf{MainStudentWindow}\\\scriptsize the hub opens};

  \node[box, right=34mm of b, fill=warnred!12, draw=warnred] (x1) {\scriptsize "Wrong credentials"};
  \node[box, right=34mm of c, fill=warnred!12, draw=warnred] (x2) {\scriptsize "Wrong credentials"};
  \node[box, right=34mm of h, fill=warnred!12, draw=warnred] (x3) {\scriptsize "Wrong code! Please try again."\\\scriptsize clear the box, \textbf{stay open}};

  \draw[flow] (a) -- (b);
  \draw[flow] (b) -- node[lbl,left]{no} (c);
  \draw[flow] (c) -- node[lbl,left]{yes} (d);
  \draw[flow] (d) -- (e);
  \draw[flow] (e) -- node[lbl,left,pos=0.4]{no} (f);
  \draw[flow] (e) -- node[lbl,right,pos=0.4]{yes} (g);
  \draw[flow] (f) |- (h);
  \draw[flow] (g) |- (h);
  \draw[flow] (h) -- node[lbl,left]{yes} (i);
  \draw[flow] (b) -- node[lbl,above]{yes} (x1);
  \draw[flow] (c) -- node[lbl,above]{no} (x2);
  \draw[flow] (h) -- node[lbl,above]{no} (x3);
\end{tikzpicture}}
\caption{Login and multi-factor verification, including the no-credentials fallback that
lets anyone who clones the repository still log in.}
\label{fig:login}
\end{figure}

\begin{jargon}{MFA, SMTP, and an App Password}
\textbf{MFA} (Multi-Factor Authentication) means proving who you are with more than one
kind of evidence: something you know (the password) plus something you have (access to the
mailbox the code was sent to). \textbf{SMTP} (Simple Mail Transfer Protocol) is the
protocol for sending mail; Python's \code{smtplib} speaks it. Gmail will not accept your
normal password from a script, so you generate an \textbf{App Password}, a separate
16-character secret scoped to one application, which is what
\code{ATTENDANCE\_SENDER\_PASSWORD} is meant to hold.
\end{jargon}

\begin{keyidea}{The best decision in the project: gate the network call, then tell the truth}
Anyone cloning this repository has no Gmail App Password. A blind \code{smtplib.SMTP(...)}
would hang on a connect or throw, and the login would be unusable. So both email paths
check first:
\begin{lstlisting}[language=Python]
sender_email = os.environ.get("ATTENDANCE_SENDER_EMAIL", "")
sender_password = os.environ.get("ATTENDANCE_SENDER_PASSWORD", "")
if not sender_email or not sender_password:
    ...
\end{lstlisting}
and then each picks a \emph{truthful} fallback for its own context. \file{verification.py}
prints the code to the terminal and says so in a dialog, so login completes.
\file{attendance.py} returns \code{False} and the caller says ``confirmation emails could
not be sent'' rather than claiming they were.

That last distinction is the whole point. The lazy version returns \code{True} and shows
``emails sent!''. This code refuses to lie to the user about something it did not do. Two
env vars drive both paths and \file{.env.example} documents them.
\end{keyidea}

\begin{honest}{Real weaknesses in the verification flow, stated plainly}
\begin{itemize}
  \item \textbf{The code is held in memory on the client and compared there.}
        \code{self.verification\_code} is a string attribute on the window;
        \code{handle\_verification} compares the typed text to it. In a desktop app the
        user owns the process, so this is theatre against a determined user. It is a
        faithful demonstration of the MFA \emph{flow}, not a security control. Say that
        rather than overselling it.
  \item \textbf{\code{random.randint} is not cryptographic.} \code{random} is a Mersenne
        Twister seeded predictably; \code{secrets.randbelow} is the right call for
        anything that is supposed to be unguessable. Given the previous point, it changes
        nothing here, but it is the correct answer to ``how would you generate that code''.
  \item \textbf{No expiry, no attempt limit.} The code lives as long as the window and can
        be retried without limit. A real MFA has a TTL and a lockout.
  \item \textbf{The env vars are misnamed for one of their two jobs.} They are called
        \code{ATTENDANCE\_SENDER\_*} and are also used for login verification, which is
        not attendance. The README and \file{.env.example} are consistent about it, so it
        is a naming debt rather than a trap, but a reviewer will ask.
  \item \textbf{Failure to send closes the window.} In \code{send\_verification\_email},
        the \code{except} shows an error and calls \code{self.close()}. But
        \code{VerificationForm} was constructed with \code{super().\_\_init\_\_(login\_parent)},
        and the login window was hidden by \code{handle\_login}. Reading the code, a send
        failure leaves the verification window closed and the login window still hidden,
        with no path back. \textbf{Inferred}, not verified: the app could not be run to
        confirm what the user sees at that moment.
\end{itemize}
\end{honest}

\begin{honest}{The show/hide password button is inverted}
\begin{lstlisting}[language=Python]
self.password_visible = False
self.toggle_password_visibility()      # called once at construction

def toggle_password_visibility(self):
    self.password_visible = not self.password_visible
    if self.password_visible:
        self.show_hide_button.setText("Show")
        self.password_textbox.setEchoMode(QLineEdit.EchoMode.Password)   # HIDDEN
    else:
        self.show_hide_button.setText("Hide")
        self.password_textbox.setEchoMode(QLineEdit.EchoMode.Normal)     # VISIBLE
\end{lstlisting}
The branch where \code{password\_visible} is \code{True} is the branch that \emph{hides}
the password. The flag means the opposite of its name. The labels happen to read correctly
if you interpret them as ``click me to Show'', and the constructor's warm-up call happens
to leave the box masked at startup, so the user-visible behaviour is right. The code says
the opposite of what it does. Rename the flag to \code{password\_hidden} and it reads
correctly.
\end{honest}

% =====================================================================
\section{The remaining screens}

\subsection{Attendance}
\label{sec:attendance}

\begin{honest}{Attendance has no table, and nothing it shows is real}
\code{setup\_attendance\_table} contains this, and this is the entire data source:
\begin{lstlisting}[language=Python]
dummy_data = [
    ("2025-03-01", "Monday", "Present"),
    ("2025-03-02", "Tuesday", "Absent"),
    ("2025-03-03", "Wednesday", "Present"),
    ("2025-03-04", "Thursday", "Absent")
]
\end{lstlisting}
Four hardcoded rows, four fixed dates in March 2025, no student id involved. There is no
\code{Attendance} model and no attendance table in the schema; verified against
\file{student.db}, which has exactly five tables and none of them is attendance.

The consequence is that marking attendance does not persist. Reopen the screen and you are
back to Present/Absent/Present/Absent regardless of what you set. The README states this
in ``What I'd do differently'' and it is the right call to state it.

What \emph{does} persist is the side effect: every press of Mark Attendance writes one
\code{Activity} row per table row, unconditionally, whether or not the value changed. So
four rows per press, forever. The shipped log's 116 rows include several of these. The
docstring on \code{handle\_mark\_attendance} even says ``check if the user changed the
status from the original dummy data'', and then does not: it re-sends and re-logs all four
every time.
\end{honest}

The email handling here is the good part, and it is the mirror image of the verification
fallback: \code{send\_attendance\_email} returns a boolean, \code{handle\_mark\_attendance}
accumulates \code{all\_emails\_sent} across the rows, and the final dialog says either
``Attendance marked and emails sent!'' or ``Attendance marked, but confirmation emails
could not be sent'' with the env-var names in the message. That is an honest status, and
it is the pattern \file{enroll.py} should have copied.

\subsection{View schedule and the Excel export}

Two modes on one table, switched by \code{self.filter\_mode}:

\begin{itemize}
  \item \textbf{Day} (the default): shows column 0, ``Current Day'', and includes only the
        courses that meet today. The filter is
        \code{if current\_day[:2] not in course\_days\_list: continue}, that is, take the
        first two characters of \code{datetime.now().strftime("\%A")} (so ``Monday''
        becomes ``Mo'') and look for it among the section's two-character day codes. This
        is the same encoding as the clash detector, and it is why the codes must be exactly
        two characters.
  \item \textbf{Week}: hides column 0 and lists every enrolled course once.
\end{itemize}

\begin{honest}{Day mode is fragile in a way that only shows up on the wrong day}
\code{current\_day[:2]} maps Monday to ``Mo'', Tuesday to ``Tu'', Wednesday to ``We'',
Thursday to ``Th'', Friday to ``Fr'', Saturday to ``Sa'', Sunday to ``Su''. That works for
this catalogue. It relies on the process locale producing English day names: on a machine
with a non-English locale, \code{strftime("\%A")} returns a translated name and the day
view silently shows nothing at all. No error, just an empty table. It also means the
default view is empty on any day the student has no classes, which looks identical to a
failure.
\end{honest}

The export is genuinely thoughtful:

\begin{lstlisting}[language=Python,caption={\file{view\_schedule.py}: do not clobber an existing export.}]
def save_data(self, excel_file_path, df):
    while not is_file_found:
        try:
            pandas.read_excel(excel_file_path)      # if this succeeds, it exists
            counter += 1
            excel_file_path = ... + f'/{name} ({counter}).xlsx'
        except FileNotFoundError:
            df.to_excel(excel_file_path, index=False)
            is_file_found = True
\end{lstlisting}

\begin{jargon}{pandas, DataFrame, openpyxl}
\textbf{pandas} is the standard Python library for tabular data. Its central object is the
\textbf{DataFrame}, a table with named columns. \code{DataFrame.to\_excel} writes it out as
a real \file{.xlsx} spreadsheet, delegating the file format to \textbf{openpyxl}, which is
why both are in \file{requirements.txt} even though no line of code imports openpyxl
directly.
\end{jargon}

\begin{honest}{Detecting a file by trying to parse it}
Probing with \code{read\_excel} and catching \code{FileNotFoundError} works, and the
auto-increment to \code{schedule\_by\_week (1).xlsx} is a nice touch that the README is
right to call out. But it reads and fully parses the whole spreadsheet just to learn
whether it exists, when \code{os.path.exists(path)} answers the same question without
opening anything. Worse, it only catches \code{FileNotFoundError}: if the existing file is
corrupt, or locked because the user has it open in Excel, \code{read\_excel} raises
something else, which is unhandled, and the export dies with a traceback in a terminal
nobody is reading.
\end{honest}

\begin{honest}{Day-mode export is disabled, and the dead check gives away why}
\code{export\_schedule} begins with \code{if self.filter\_mode == "Day":} and returns
``Export Not Available''. Then, forty lines later, it asks \code{if self.filter\_mode ==
"Day":} again, and that second block is unreachable. The README explains the history: the
day view hides a column, the export was writing misaligned data, so it was disabled rather
than shipped broken. That is the right decision. The residue is the dead second check and a
\code{columns} list assigned in the first branch and never used.

The consequence worth knowing for a demo: \code{filter\_mode} starts at \code{"Day"}, so
\textbf{Export does nothing until you press Week first}. That looks like a broken button.
\end{honest}

\begin{honest}{The schedule window logs an activity in its constructor}
\code{\_\_init\_\_} writes a \code{"Filter by Day"} activity row before the user has done
anything, simply because the screen opened in day mode. So the log conflates ``opened the
schedule'' with ``pressed the Day button''. Combined with the attendance screen's four
rows per press, this is why the shipped log has 116 rows.
\end{honest}

\subsection{Activity log}

The simplest screen: read every \code{Activity} row for the student, render six columns.
It is the only place the log is consumed.

\begin{honest}{A guard that hides data}
\begin{lstlisting}[language=Python]
if activity_log_data.course_id is not None and activity_log_data.course_section is not None:
    self.activity_label_table_widget.setItem(row, 2, QTableWidgetItem(str(...course_id)))
self.activity_label_table_widget.setItem(row, 3, QTableWidgetItem(...course_section))
\end{lstlisting}
Column 2 (Course ID) is populated only if \emph{both} the id and the section are non-null.
So a log entry with a course id but no section shows a blank Course ID column, even though
the id is right there in the row. Column 3 has no such guard and is set unconditionally,
which for a null section passes \code{None} into \code{QTableWidgetItem}. There is no
sorting and no pagination: 116 rows today, and it grows without limit because nothing ever
prunes it.
\end{honest}

\subsection{The window-transition machinery}

Every screen defines \code{closed = pyqtSignal()} and emits it from \code{closeEvent}. The
hub connects that signal to its own \code{reopen\_login} (which just calls
\code{self.show()}), and hides itself when opening a child. That is how closing a child
brings the hub back.

\begin{jargon}{Signals and slots}
Qt's event mechanism. An object \textbf{emits} a \textbf{signal} (``I was closed'') without
knowing or caring who is listening. Any number of functions, called \textbf{slots}, can be
\code{connect}ed to it and are called when it fires. It is Qt's version of the observer
pattern, and it is what decouples the child windows from the hub.
\end{jargon}

\begin{honest}{\code{closeEvent} never calls \code{super()}, and the hub calls it by hand}
Every \code{closeEvent} override in the project has \code{\# super().closeEvent(event)}
commented out and never accepts or ignores the event. Qt's default is to accept, so the
windows do close, but the override is doing something other than what its signature
promises.

\code{MainStudentWindow.logout} is stranger still: it sets a flag and then calls
\code{self.closeEvent(event=None)} \emph{directly}, invoking an event handler as an
ordinary method with a fabricated null event. It works only because the body never touches
\code{event}. The intended API is \code{self.close()}, which lets Qt deliver a real event.

There is also a straightforward duplicate in \code{MainStudentWindow.activity\_log}:
\begin{lstlisting}[language=Python]
self.activity_log_window = ActivityLogWindow()
self.activity_log_window = ActivityLogWindow(email=self.email)
\end{lstlisting}
Two windows are constructed; the first is immediately orphaned by the reassignment. It is
never shown, but its constructor ran, which means it opened the \file{.ui} file, ran
\code{get\_std\_id\_by\_email} for the \emph{default} student (Alice), and queried her
activity log. A wasted window and a wasted round of queries on every click.
\end{honest}

\subsection{The name collision}

\begin{honest}{Two different classes are both called \code{ViewScheduleWindow}}
\file{enroll.py} defines \code{class ViewScheduleWindow(QMainWindow)}, which loads
\file{enroll.ui} and is the \emph{checkout} screen. \file{view\_schedule.py} defines
\code{class ViewScheduleWindow(QMainWindow)}, which loads \file{view\_schedule.ui} and is
the \emph{schedule} screen. Same name, different files, different purposes.

Nothing breaks, because \file{main\_window.py} imports the schedule one and
\file{add\_courses.py} imports the enroll one, and no module imports both. It is a trap
lying in wait for the first person who needs both in one file, and it makes the codebase
hard to talk about. The history is visible: \file{enroll.py} is the older file (its class
was clearly copied from the schedule screen and repurposed), and the \file{.ui} it loads is
still named \code{view\_schedule\_window} internally, as is \file{add\_courses.ui}. Three
files carry the fingerprint of one original. The README's final ``Challenges'' bullet
already records this collision, which is the right instinct.
\end{honest}

\begin{honest}{An orphan interface file}
\file{add\_remove\_interaction.ui} is 8.5 KB, the largest \file{.ui} in the project, and no
Python file loads it. Verified by grep. It is a designed screen that was never wired up.
\end{honest}

% =====================================================================
\section{Packaging and running}
\label{sec:docker}

\begin{jargon}{Docker, image, container}
\textbf{Docker} packages an application together with its operating-system dependencies
into an \textbf{image}, a frozen filesystem plus a start command. Running an image gives
you a \textbf{container}, an isolated process that behaves the same on any machine with
Docker. It is the standard answer to ``it works on my machine''.
\end{jargon}

The \file{Dockerfile} is short and its shape is right: start from \code{python:3.10-slim},
install the system libraries Qt needs, set \code{WORKDIR /app} (which is what makes the
relative \code{loadUi} paths resolve), install \file{requirements.txt} \emph{before}
copying the code so that Docker's layer cache is not invalidated by every edit, copy
everything including the seeded \file{student.db}, set \code{QT\_QPA\_PLATFORM=xcb}, and
run \file{login\_page.py}.

\begin{jargon}{X11 forwarding, and why a GUI in a container is awkward}
A container has no screen. On Linux the display is managed by a server called \textbf{X11}
which listens on a socket at \file{/tmp/.X11-unix}. The run scripts mount that socket into
the container and pass the \code{DISPLAY} variable, so the containerised app draws on
\emph{your} desktop. \code{xhost +local:docker} is the command that permits it, and it is a
blunt instrument: it grants any local container access to your display for the session.
\end{jargon}

\begin{honest}{Docker can never send an email, by construction}
Verified by grep: neither the \file{Dockerfile} nor any of the three run scripts mentions
\code{ATTENDANCE\_SENDER\_EMAIL} or \code{ATTENDANCE\_SENDER\_PASSWORD}. No \code{ENV}, no
\code{-e} flag, no \code{--env-file}. A container gets a fresh environment, so those
variables are always empty inside it.

The app therefore \emph{always} takes the fallback path in Docker: the verification code is
printed to the container's terminal and attendance always reports that emails could not be
sent. Because the fallback is honest, this degrades gracefully instead of breaking, which
is a real vindication of the design. But the README's claim that the scripts ``forward the
host display into the container'' is the whole story only for the display. If you demo from
Docker, email is off, and adding
\code{-e ATTENDANCE\_SENDER\_EMAIL -e ATTENDANCE\_SENDER\_PASSWORD} to the \code{docker
run} lines is a one-line fix per script.
\end{honest}

\begin{honest}{The Dockerfile probably does not build any more. Labelled inferred.}
The apt line installs \code{libgl1-mesa-glx}. The \code{python:3.10-slim} tag today is
built on Debian 12 (bookworm), and \code{libgl1-mesa-glx} was \textbf{removed} in bookworm,
replaced by \code{libgl1} and \code{libglx-mesa0}. If that is so, \code{apt-get install}
exits non-zero and the image build fails at step 2.

\textbf{This is inferred and was not verified.} The Docker daemon was not running on the
machine this document was written on, so no build was attempted. Treat it as the first
thing to check before demoing from Docker: run \code{docker build -t pscs-app .} and see.
If it fails, change \code{libgl1-mesa-glx} to \code{libgl1} and it should proceed. Note
this is a bit-rot issue, not a mistake: the line was almost certainly correct against the
base image of the day.
\end{honest}

\begin{honest}{The run scripts contain characters that are not plain ASCII}
All three (\file{run-linux.sh}, \file{run-macos.sh}, \file{run-windows.ps1}) contain a
hammer emoji in their build message and a non-breaking hyphen (Unicode U+2011, not the
ordinary U+002D) in the phrase ``Auto-build'' in their header comment. Both are in
comments and echo strings, so they are harmless, but they are the fingerprint of text
pasted from a rich-text editor. The non-breaking hyphen in particular is invisible in most
editors and is the sort of thing that breaks a naive grep. (Those exact bytes are described
here rather than reproduced, because this document is typeset with a font stack that has no
glyph for either one.)

\file{queries.py} and \file{verification.py} carry the same fingerprint: smart apostrophes
inside docstrings, where an ASCII apostrophe was meant.
\end{honest}

\subsection{Dependencies, and what each is for}

\begin{center}
\begin{tabular}{@{}lll@{}}
\toprule
\textbf{package} & \textbf{pinned} & \textbf{role} \\
\midrule
PyQt6 & \code{>=6.6,<7} & every window, widget, dialog and signal \\
SQLAlchemy & \code{>=2.0,<3} & the ORM; models and every query \\
pandas & \code{>=2.0} & builds the DataFrame for the Excel export \\
openpyxl & \code{>=3.1} & the \file{.xlsx} writer pandas delegates to \\
\bottomrule
\end{tabular}
\end{center}

The upper bounds on the first two are correct practice: a major version bump of either is
a breaking change. \code{pandas} and \code{openpyxl} have no upper bound, which is a minor
inconsistency. The whole dependency set exists to serve one feature each, and there is
nothing gratuitous in it.

% =====================================================================
\section{The verification story, and what it does not cover}

\begin{honest}{The README describes a test methodology that no file in the repo implements}
The README's Challenges and ``What I learned'' describe, in specific detail, a headless
smoke test: setting \code{QT\_QPA\_PLATFORM=offscreen}, monkeypatching the static
\code{QMessageBox} methods (\code{information}, \code{warning}, \code{critical}) so modals
record their calls instead of blocking, and driving login through to the main window plus
every sub-screen in one pass. It says ``This is how the project was last verified.''

There is \textbf{no test file in the repository}. \code{git ls-files} lists no
\file{tests/} directory, no \file{test\_*.py}, no pytest configuration, and no CI workflow.
\file{.gitignore} mentions \code{.pytest\_cache/}, which is boilerplate.

The most likely reading, and this is \textbf{inferred}, is that the smoke test was written
as a throwaway script during a verification pass and never committed. That is a perfectly
honest thing to have done, but the README's present tense reads as though the harness is
part of the project, and an interviewer who goes looking for it will not find it. Two
options, both fine: commit the script, or reword the README to say the test was run once,
ad hoc, and not kept. The README's own last bullet already concedes ``Add tests. Everything
here was verified by clicking through the app'', which sits oddly beside the detailed
description of the harness. The two passages should be reconciled.
\end{honest}

\begin{keyidea}{What this document verified, and what remains open}
\textbf{Verified by execution:} the capacity arithmetic and its overflow past
\code{max\_capacity}; the swap's data loss; the foreign keys' invalidity under
\code{PRAGMA foreign\_keys = ON}; \code{bool(Qt.CheckState.Unchecked) == True}; the
\code{glob\_id} default of 20180; the ORM relationship warning. \textbf{Verified by
inspection of \file{student.db}:} every row count, the day encodings, the over-capacity
rows, the polluted \code{activity.course\_id} values, the all-NULL dead columns, the empty
\code{filter} table. \textbf{Verified by grep:} \code{max\_capacity} unread,
\file{courses\_database.py} unimported, \file{add\_remove\_interaction.ui} unloaded,
\code{insert\_attendance\_activity} uncalled, the env vars unforwarded by Docker.

\textbf{Open, and labelled inferred where claimed:} anything about on-screen behaviour
(PyQt6 was not present, and the app needs a desktop); the Dockerfile's
\code{libgl1-mesa-glx} problem; the verification window's dead end on a send failure; the
history behind the missing test file.
\end{keyidea}

% =====================================================================
\section{Findings, ranked}

If you have five minutes before an interview, read only this.

\begin{enumerate}
  \item \textbf{The swap loses data} (\S\ref{sec:swap}). Delete commits, then insert is
        attempted; a failure between them destroys the enrollment and the UI says nothing
        happened. Verified by running it. Fix: one transaction around both.
  \item \textbf{The foreign keys are structurally invalid and never enforced}
        (\S\ref{sec:fkbroken}). They reference one column of a composite primary key.
        SQLite's default of not checking is the only reason the app works. Verified.
  \item \textbf{\code{max\_capacity} is never read} (\S\ref{sec:capacity}), and the shipped
        database already has two sections at 51 seats against a maximum of 50. Verified.
  \item \textbf{Passwords are plaintext} and compared with \code{==}. Already documented in
        the README, which is the right way to carry it.
  \item \textbf{The clash rule rejects back-to-back classes} (\S\ref{sec:clash}) because
        all three interval tests use \code{<=} on the boundary.
  \item \textbf{Attendance is four hardcoded rows} with no table behind them
        (\S\ref{sec:attendance}), and it writes four log rows per press regardless of
        change.
  \item \textbf{The README is wrong in three places}: \file{courses\_database.py} is not a
        seeding helper (it is a dead, incompatible \code{Student} model); \code{glob\_id}
        is per-instance so the stated call-order hazard does not exist as described (the
        shared \emph{session} is the real one); and the described headless test harness is
        not in the repository.
  \item \textbf{Dead weight:} the \code{filter} table (0 rows, no code), the \code{swap\_id}
        and \code{remove\_id} columns (all NULL), \file{courses\_database.py},
        \file{add\_remove\_interaction.ui}, \code{insert\_attendance\_activity}, the
        double-construction in \code{activity\_log()}, and the unreachable second day check
        in \code{export\_schedule}.
  \item \textbf{\code{python queries.py} wipes the \code{activity} and \code{course}
        tables} with no confirmation.
  \item \textbf{Cosmetic but demo-visible:} the drop screen's headers are off by one from
        column 6; Export appears broken until Week is pressed; ``Test Course'' at 02:30 on
        MoSa is in the catalogue.
\end{enumerate}

\begin{keyidea}{And what is genuinely good, which is worth saying too}
The \code{AllQueries} choke point is a real architectural decision that paid off, and the
README's account of why (a two-person team needed one contract between screens) is
credible and supported by the code. The cart/enroll split is the right model. The
environment-gated email with a \emph{truthful} fallback in both paths is better engineering
than most coursework contains, and it is the reason the app still runs for anyone who
clones it. The Excel auto-increment shows someone thought about the second time the button
is pressed. And disabling the day export rather than shipping a misaligned spreadsheet was
the right call, made for the right reason.

Above all: the README already names most of the weaknesses in this document, in public, on
a portfolio piece. That is the hardest thing on this list to do, and it is the thing to
lead with.
\end{keyidea}

\end{document}
